% results  -  part of main.tex (see journal/README.md)
\section*{Results}

\subsection*{A memory block addressed by the token itself}

The block has three parts (Fig.~\ref{fig:arch}a, Methods and Supplementary Note~1). A \emph{key} is the
token's layer-0 embedding without position, and the read head's query projection is
initialised to the identity and frozen, so a token matches itself exactly. A \emph{read} scores the stored
keys by cosine similarity, sharpens the scores by an annealed softmax in the manner of a modern Hopfield network\cite{hopfield1982,ramsauer2021hopfield} (an associative memory that returns the stored pattern closest to a cue), and injects the blended values into the residual stream (the running representation that every layer reads and updates) through a gate initialised at zero\cite{bachlechner2021rezero}.
At initialisation the block is therefore an exact no-op, and the store-disabled model is a
true control. A \emph{write} stores one slot per token, with the value shifted by one
position so that the slot keyed on a token returns what followed it.

Where reads and writes happen matters. Writing the mean of all layers' states into one slot
per token gives 80.8\% recall on a 32-fact task. Letting every layer write its own
post-block state as its own slot in the one shared store lifts this to 95.5\% and 98.9\% over two seeds (Table~\ref{tab:core}; Supplementary Note~2 and Supplementary Table~1). A control with the same
number of slots, but keyed by each layer's evolving residual, stays at
0.2\%. The gain comes from per-depth values that stay exactly addressable. Extra capacity
on its own does not do it. The exact address is load-bearing. Every attempt to enrich the
key---moving it deeper, adding position, mixing in a neighbouring token, or projecting it to
fewer dimensions---collapsed matching to chance (Supplementary Note~2).

\subsection*{Recall of facts stated once, out of context}

We test memory with a task in which no answer can come from the prompt (Supplementary Note~3). The first segment
states $n$ key--value facts inside PG-19 prose (a corpus of books)\cite{rae2020compressive}; intermediate
segments are unrelated prose; the final segment holds only a query marker and one key. Two
matched controls run every time: the store disabled, and a scrambled store in which each
example reads a neighbour's store.

With single-token sub-word keys\cite{sennrich2016bpe}, recall is 100\% at 8, 16 and 32
facts and 99.1\% at 64, with the scrambled control at $0.16\%$ (Table~\ref{tab:core}; Supplementary Table~2).
Memory survives the process: a store written, then reloaded in a fresh interpreter, answers
at 100\%, and a mismatched store file at $0.0\%$. Two obstacles had to be removed on
the way. Character-level keys impose a \emph{multi-token-key tax}: the only token at the
retrieval point is a shared delimiter, so recall falls from
90\% to 55\% between 4 and 13 facts; a single-token key removes it. The
familiar fact-count cliff---recall collapsing beyond a threshold even though the top-scoring
slot is correct---is a training-dynamics deadlock rather than a capacity limit. Multi-layer
placement clears it from a cold start after a delayed phase transition (99.7\% at 16
facts, 98.0\% at 32), and a warm-start curriculum reaches the same endpoint in about
half the steps.

\subsection*{The weights are stateless; the store is the state}

Splitting weights from store predicts that a store can be moved between conversations
without touching the model. We hold one trained model fixed and build two stores from
disjoint fact sets (Fig.~\ref{fig:arch}b and Supplementary Note~4). Recall
is diagonal: 100\% and 96.9\% when the store that holds the queried facts is
loaded, and $0.0\%$ when the other store is, with a scrambled store at chance (one of 32 queries, $3.1\%$; Table~\ref{tab:core} and Supplementary Table~3). Swapping the store swaps the memory.

The same experiment marks the boundary. Two independently seeded models each read their own
store at 98.8\% and 100\%, but each other's at $0.0\%$, in both directions,
although the facts are physically present and the stores are shape-compatible. A store's
keys are its writer's token embeddings, so a store is a memory only for the model that wrote
it, and the persistence path refuses a foreign store by its model fingerprint. One frozen
model can therefore serve many users by loading their stores. Carrying memory across model
versions requires re-encoding, which is cheap: the stored keys decode exactly to their
tokens, so the key sequence is itself the transcript of what was written (Methods).

\subsection*{Memory far beyond the attention window}

Because the store sits outside attention, it can span many attention windows. A fixed ring
buffer eventually overwrites old facts, which is a limit of eviction policy and not of
capacity. We
instead grow the store by appending, like a key--value cache, and read it with a top-1 lookup, which uses only the best-matching slot (Supplementary Note~5). Recall then stays at 100\% from 1.9k to 998k tokens between fact and query (Fig.~\ref{fig:long}a and Supplementary Table~4), about 7,800 times the 128-token window, on both a
two-layer and a six-layer model. A dense read over all slots decays as distractors
accumulate, from 85\% to 60\% by 25k tokens and to $0\%$ by 128k, so sparse reading is a
requirement, not a compromise. A scrambled store sits at $0.0\%$ at 128k and 384k tokens.
Reaching a million tokens needed an evaluation-time change only: building the store in place
rather than rewriting it every segment (Methods). The read is $O(1)$ per query, so build time
sets the reach. One caveat runs the other way. Keys drawn from
the rarest words, whose embeddings are least trained, recall at 50--75\% even at a million
tokens. That is an effect of embedding quality, not of store size.

The store also behaves as a shared key--value memory in the sense used to compress attention caches\cite{shazeer2019mqa,ainslie2023gqa,deepseekv2mla2024,deepseekv41flash2026} (Supplementary Note~6). An index of
each slot's key token lets a query score about 2 of 8,001 slots at the recall of a dense scan
(81.9\% versus 82.5\%; Supplementary Table~5). Building that candidate pool once and reusing it at every read layer
gives recall identical to rebuilding it per layer. The stored codes survive quantization (storing numbers at reduced precision) to 4 bits per element with no loss (Fig.~\ref{fig:long}b), eight times smaller than the
32-bit store. Keys are the
bit-limiting half: they fail at 3 bits, while values hold to 3. Both figures are measured
properties of the read. Turning them into speed and memory savings needs specialised low-level routines for hashed lookup and 4-bit integers.

\subsection*{No lost-in-the-middle, but a write-order limit}

Attention models retrieve facts placed in the middle of a long context worse than those at
its ends\cite{liu2024lost}. We tested this with a passkey design\cite{mohtashami2023landmark}, in which a fact is planted at a chosen depth in long filler text and queried at the end, varying that depth (Supplementary Note~5). The store shows no mid-context dip:
recall is monotonic in depth (Fig.~\ref{fig:passkey}a). In its place there is a \emph{write-order}
decay that grows with saturation. A fact written first stays at 100\%, but a fact
written last falls from 81\% to 50\% to 6\% as the store fills from 4k to 16k to 65k slots,
with a scrambled store at $0\%$ throughout. The store privileges what it saw first, the
opposite of a recency bias.

Two ablations locate the cause (Fig.~\ref{fig:passkey}b). Ring overwriting is not the
explanation: an append-only store gives an identical curve. The write policy is part of it. The block stores a
segment only when it is novel enough relative to the store (Methods), and a late fact whose
key token recurs in the intervening prose reads as familiar. Keying facts on unique tokens,
which always read as novel, lifts the late tail fourfold, from 6\% to 25\%, though it lowers
the early end from 100\% to 75\% because rarely seen tokens have weaker embeddings. On a
weaker six-layer model the decay is nearly binary: only the first-written fact survives,
and unique keys do not help (Fig.~\ref{fig:passkey}c). The limit is part fixable write
policy, part saturation, and nothing to do with position.

\subsection*{From a lookup table to a memory}

A store that can only be appended to and read is a lookup table. A memory must also revise a
fact when it is restated, stay bounded as text accumulates, and support computation over its
contents. We added each by changing only the read and write heads (Fig.~\ref{fig:memory}; Supplementary Note~7 and Supplementary Table~6).

\emph{Revision.} An append-only store keeps both the old and new value of a restated fact and
cannot tell which is current. A small bias towards the most recently written match, read from
a global write counter, lifts newest-value recall from 88.9\% to 96.7\%, while facts that were never updated stay at 97.5\% (Fig.~\ref{fig:memory}a). The
position-free key carries no recency of its own; recency has to come from write order.

\emph{Boundedness.} Gated write operations keyed on the token---\textsc{new},
\textsc{update-in-place} or \textsc{skip}---keep one slot per distinct fact. Recall is
95.9\% at about 200 occupied slots against 94.1\% at 508 when every occurrence
is appended, and updates land at 96.2\% against 7.8\%
(Fig.~\ref{fig:memory}b). A single shared store with these operations matches separate
per-layer stores (96.2\% recall, 96.6\% update).

\emph{Composition.} Given $A\!\to\!B$ and $B\!\to\!C$ stored beyond the window, answering $A$
with $C$ needs two retrievals chained through the store. Feeding a retrieved representation
back as the next query failed in all seven variants we tried, at chance: a stored value does
not lie on the manifold of keys, so the next exact match never fires. Composition works
as a working-memory operation. The model emits each intermediate as a token and
queries again, a chain of thought\cite{wei2022cot} (writing out intermediate steps) whose steps are store reads. Joint recall
is 96.9\% at two hops and 95.2\% at three over 64 facts, and 91.6\% at
four and 89.7\% at five over 32 facts, close to 0.98 per hop (Fig.~\ref{fig:memory}c). Feeding the model its own emitted intermediates, in place of the true
ones, changes nothing. Breaking one link of the chain drops joint recall to $0.0\%$ while
the preceding hops still succeed, so the model is composing rather than taking a shortcut.

\emph{Meaning.} The exact match is not semantic. A synonym of a stored key retrieves at
chance (about 4\%), and softening the read narrows what matches further. Meaning
can come from reasoning instead. When the first hop maps a synonym to its canonical form from
the weights and the second hop reads the store, synonym-queried facts are recalled at
90.8\%; replacing the synonym map with a random, unlearnable one drops the first hop
to 2\%. The store is an exact index, and meaning is supplied above it.

